% Generated draft: requires mathematical and source review.
\documentclass[11pt]{article}
\usepackage[margin=25mm]{geometry}
\usepackage{fontspec}
\setmainfont{Latin Modern Roman}
\usepackage{amsmath,amssymb,mathtools,unicode-math}
\setmathfont{Latin Modern Math}
\usepackage{enumitem,needspace,xurl,hyperref}
\setlength{\parindent}{0pt}
\setlength{\parskip}{5pt}
\setlength{\emergencystretch}{3em}
\newcommand{\N}{\mathbb{N}}
\newcommand{\Z}{\mathbb{Z}}
\newcommand{\Q}{\mathbb{Q}}
\newcommand{\R}{\mathbb{R}}
\newcommand{\C}{\mathbb{C}}
\newcommand{\F}{\mathbb{F}}
\newcommand{\A}{\mathbb{A}}
\newcommand{\PP}{\mathbb{P}}
\newcommand{\E}{\mathbb{E}}
\newcommand{\eps}{\varepsilon}
\newcommand{\dd}{\,\mathrm{d}}
\begin{document}
\section*{UnsolvedMath ID: 20002890}
\textbf{Source problem code:} catalogue entry
\subsection{Definitions and mathematical statement}
Han Peters's source problem takes a rational proper map $R:\mathbb B^n\to\mathbb B^N$, $n\ge3$, of degree $d=d_0(R)$ and asks: is
\[
 d\le\frac{N-1}{n-1}?
\]
Does the same bound hold under the additional assumption $n\ge2d^2+2d$, and which maps attain equality? It records an affirmative monomial result for the latter question and, for $n=2$, the sharp monomial bound $d\le2N-3$. A further question asks, separately in the rational, polynomial, or monomial class, whether any degree $d>1$ occurring at fixed $(n,N)$ implies existence of a map in the same class of degree $d-1$. This asks for another map, not a deformation of the original map.

Use the usual rational degree $\max\{\deg p_i,\deg q\}$ after common factors are removed from $R=p/q$; the source does not otherwise define $d_0$. A monomial map has coordinates constant multiples of holomorphic monomials. The aid explicitly allows zero target coordinates, so the map need not be linearly full.

For $N\ge n$ its reported monomial degree spectrum is
\[
 \{d_0(f):f:\mathbb B^n\to\mathbb B^N\text{ proper monomial}\}
 =\begin{cases}
 \{1,\ldots,2N-3\},&n=2,\\
 \{1,\ldots,\lfloor(N-1)/(n-1)\rfloor\},&n\ge3.
 \end{cases}
\]
The aid attributes the upper bounds to D'Angelo--Kos--Riehl and Lebl--Peters, and combines them with an explicit generalized Whitney construction and zero-coordinate padding to realize every integer in the intervals. The full construction and precise published bound hypotheses are not supplied here. Under that convention the degree-lowering question is affirmative for all proper monomial maps, at every occurring $d>1$.

Let $n\ge3$ and $D=\lfloor(N-1)/(n-1)\rfloor$. If a rational or polynomial map fails degree lowering, then
\[
 d\ge D+2.
\]
Indeed if $d-1\le D$, a monomial map supplies the requested lower-degree example in either broader class. This is a necessary obstruction to failure, not an existence theorem for counterexamples at higher degree. The interval theorem does not prove the full rational degree bound, settle its large-$n$ variant for all rational maps, or classify every equality map. Requiring linearly full targets would be a separate version not covered by the stated padding argument.
\subsection{Short English statement}
The aid reports exact monomial degree intervals when zero target coordinates are allowed, which gives monomial degree lowering. Any rational or polynomial failure must therefore occur at least two integers above the monomial endpoint. The general rational bounds and classification of equality maps remain separate questions.
\subsection{Sources}
\begin{itemize}
\item Local UnsolvedMath registry, permanent ID 20002890, source code catalogue entry; created and updated 14 August 2026. Canonical \texttt{aim-several-complex-variables-notes.json}, zero-based \texttt{notes[47]}, tag \texttt{problem}, item 8 proposed by Han Peters.
\item AIM ``Complexity of mappings in CR geometry'': \url{https://aimath.org/WWN/complexitycrmap/complexitycrmap.pdf}. Wording is marked exact, while supplied verification reports restoring ball superscripts, the fractions, degree notation, and a page header. Original remarks and literature field are empty.
\item The original cites [DLP] for the conditional monomial claim and [DKR] for the sharp $n=2$ bound. The aid names D'Angelo--Kos--Riehl and Lebl--Peters, but supplies no full titles, dates, theorem statements, bibliography entries, or generalized Whitney formulas.
\item Research attempt 1: \texttt{partial\_result}, problem status at run \texttt{partially\_solved}, rigor label \texttt{proved}, novelty confidence low. Ulam AI / UnsolvedMath Contributors, CC BY 4.0; machine-generated research aid, not peer reviewed, independently verify before citation.
\end{itemize}
\end{document}
